\section{Analyse des besoins}\label{sec:besoins}

\subsection{Acteurs}

\begin{itemize}
	\item \textbf{Visiteur} : consulte la disponibilité du réseau, s'inscrit, se connecte.
	\item \textbf{Usager} : souscrit et gère son abonnement, emprunte et restitue des vélos, consulte ses trajets et ses factures, signale un problème.
	\item \textbf{Opérateur} (régulation ou maintenance) : suit la flotte et les locations, traite les signalements, consulte les statistiques et les propositions de rééquilibrage.
	\item \textbf{Administrateur} : droits de l'opérateur, plus la gestion des stations, des bornes et des comptes.
\end{itemize}

\diagramme{use_cases}{Diagramme de cas d'utilisation}

Le diagramme (figure~\ref{fig:use_cases}) organise les acteurs en hiérarchie : l'\textbf{utilisateur authentifié} est l'acteur générique (s'authentifier, signaler un problème) ; l'\textbf{usager} et l'\textbf{opérateur} en héritent, et l'\textbf{administrateur} est un opérateur doté de droits supplémentaires (comptes, description des stations). Dans l'application, un visiteur non connecté peut en outre consulter la disponibilité du réseau et s'inscrire.

Les cas d'utilisation se regroupent en trois familles :
\begin{itemize}
	\item \textbf{usage du réseau} (usager) : s'abonner, gérer sa souscription, se désabonner, emprunter et rendre un vélo ;
	\item \textbf{exploitation} (opérateur) : tournées de rééquilibrage, opérations de maintenance, consultation des statistiques ;
	\item \textbf{administration} : gestion des comptes et description des stations.
\end{itemize}

Un acteur secondaire, le \textbf{système}, représente les traitements automatiques de la base. Les relations \og include \fg{} montrent qu'emprunter ou rendre un vélo déclenche toujours l'historisation de son état et de sa position, et que l'abonnement entraîne la facturation ; les relations \og extend \fg{} marquent les cas facultatifs, comme la gestion ou la résiliation d'une souscription existante. Ces inclusions correspondent directement aux déclencheurs et aux événements planifiés de la partie~\ref{sec:implementation}.

\subsection{Règles métier}

Les règles suivantes, issues du sujet, sont garanties par la base elle-même et non par l'application (partie~\ref{sec:implementation}) :
\begin{enumerate}
	\item un vélo ne peut être qu'à un seul endroit à la fois (sur une borne, en location ou en atelier) ;
	\item une borne n'accueille qu'un seul vélo ;
	\item un usager n'emprunte qu'un vélo à la fois, et seulement avec un abonnement valide ;
	\item un vélo signalé défectueux ou en atelier ne peut pas être emprunté ;
	\item la durée et le coût d'une location sont calculés automatiquement à la restitution ;
	\item tout changement d'état ou de position d'un vélo est historisé.
\end{enumerate}
